\chapter{Dataset Sources}
\label{app:datasets}

This appendix documents every dataset bundled with the book, its
format, its provenance, and how to obtain more data from public
sources for extending the examples. All bundled files live in
\texttt{crates/quant-lab/data/} and are loaded by the
\crate{quant-lib} examples via the helpers in
\texttt{examples/common/mod.rs}.

\section{Bundled Datasets}

\subsection{Synthetic Stock Prices (\texttt{stock\_prices.csv})}

A 63-row synthetic OHLCV series for a single fictitious stock from
2 January 2024 to 28 March 2024. Used by the early examples (Ch.~3, 4,
5) where a clean, controlled series is preferable to real market data.

\begin{center}
\begin{tabular}{ll}
\toprule
\textbf{Field} & \textbf{Description} \\
\midrule
\texttt{date} & Trading date (ISO 8601) \\
\texttt{open, high, low, close} & Daily OHLC prices \\
\texttt{volume} & Daily traded volume \\
\texttt{adj\_close} & Dividend/split-adjusted close (equals close here) \\
\bottomrule
\end{tabular}
\end{center}

\begin{lstlisting}[language=Rust]
use quant_lib::prelude::*;
let path = common::stock_csv_path(env!("CARGO_MANIFEST_DIR"));
let bars = common::load_stock_csv(&path);
\end{lstlisting}

\subsection{B3 Brazilian Stocks (8 JSON files)}

Daily OHLCV data for eight liquid stocks traded on B3 (Brasil Bolsa
Balcão), the S\~ao Paulo exchange. Each file contains 1{,}247 daily
bars from 22 July 2021 to 26 July 2024---roughly three years of
trading. Prices are in Brazilian Reais (BRL).

\begin{center}
\begin{tabular}{lll}
\toprule
\textbf{Symbol} & \textbf{Company} & \textbf{Sector} \\
\midrule
\texttt{PETR4} & Petr\'oleo Brasileiro (Petrobras) & Energy \\
\texttt{VALE3} & Vale & Mining \\
\texttt{ITSA4} & Ita\'usa (Ita\'u) & Financial \\
\texttt{BBDC4} & Bradesco & Financial \\
\texttt{B3SA3} & B3 (exchange operator) & Financial \\
\texttt{ABEV3} & Ambev & Consumer staples \\
\texttt{GGBR4} & Gerdau & Steel \\
\texttt{WEGE3} & WEG & Industrial \\
\bottomrule
\end{tabular}
\end{center}

\marginnote{%
\textbf{Why Brazilian stocks?}\\[0.3em]
B3 daily data is freely available, covers liquid names, and gives the
examples a consistent cross-sectional universe for portfolio and
factor exercises. The same code works with any OHLCV source---US
equities, ETFs, or crypto---by swapping the data loader.
}

\paragraph{JSON format.}
Each file is a JSON array of objects with keys
\texttt{date, open, high, low, close, volume}. The bundled parser
in \texttt{common/mod.rs} is a minimal splitter that relies on the
fixed key order; for general JSON, use \texttt{serde\_json}.

\begin{lstlisting}[language=Rust]
use quant_lib::prelude::*;
let path = common::b3_json_path(env!("CARGO_MANIFEST_DIR"), "PETR4");
let bars = common::load_json_ohlcv(&path);
\end{lstlisting}

\subsection{Credit Card Fraud Sample (\texttt{creditcard\_sample.csv})}

A 20-row anonymised sample from the Kaggle Credit Card Fraud Detection
dataset. Each row is one transaction; columns \texttt{V1}--\texttt{V28}
are PCA-transformed features, \texttt{Amount} is the transaction value,
and \texttt{Class} is $1$ for fraud, $0$ otherwise. Used in Ch.~1--2
for the fraud-detection and loan-default classifiers.

\begin{center}
\begin{tabular}{ll}
\toprule
\textbf{Field} & \textbf{Description} \\
\midrule
\texttt{Time} & Seconds elapsed since first transaction \\
\texttt{V1}--\texttt{V28} & PCA-transformed features (anonymised) \\
\texttt{Amount} & Transaction value (EUR) \\
\texttt{Class} & $1$ = fraud, $0$ = legitimate \\
\bottomrule
\end{tabular}
\end{center}

\begin{warning}
The bundled sample is too small for real modelling. It exists only to
make the chapter examples runnable offline. For actual experiments,
download the full 284{,}807-row dataset from Kaggle (see below).
\end{warning}

\section{Public Data Sources for Extending the Examples}

The bundled data is enough to run every example in the book. To go
further---longer histories, cross-sectional studies, alternative
asset classes---the following sources are all free and downloadable
without registration beyond an email.

\subsection{Equity OHLCV}

\begin{center}
\begin{tabular}{lll}
\toprule
\textbf{Source} & \textbf{URL} & \textbf{Format} \\
\midrule
Stooq & \texttt{stooq.com/q/d/?s=\{symbol\}\&i=d} & CSV OHLCV \\
Yahoo Finance & \texttt{finance.yahoo.com/quote/\{TICKER\}/history} & CSV OHLCV \\
Binance API & \texttt{api.binance.com/api/v3/klines} & JSON OHLCV \\
Alpha Vantage & \texttt{alphavantage.co} (free tier) & JSON/CSV \\
\bottomrule
\end{tabular}
\end{center}

Stooq is the easiest for bulk download: it serves daily and intraday
bars for US, European, and Brazilian tickers with no API key. Yahoo
Finance covers the widest universe. Binance is the primary source for
crypto OHLCV via a clean REST API.

\subsection{Factor and Macro Data}

\begin{center}
\begin{tabular}{lll}
\toprule
\textbf{Source} & \textbf{URL} & \textbf{Content} \\
\midrule
Ken French Data Library & \texttt{mba.tuck.dartmouth.edu/pages/faculty/ken.french} & Fama-French factors, momentum, industry portfolios \\
FRED (St.\ Louis Fed) & \texttt{fred.stlouisfed.org} & Macro series (rates, CPI, vix) \\
AQR Data Sets & \texttt{aqr.com/Insights/Datasets} & Factor returns, Betting-Against-Beta \\
World Bank & \texttt{data.worldbank.org} & Country-level macro indicators \\
\bottomrule
\end{tabular}
\end{center}

The Ken French library is the canonical source for the factor data
used in Ch.~12. FRED is the standard for US risk-free rates
(\texttt{DGS3MO}, \texttt{DGS10}) and the VIX (\texttt{VIXCLS}).

\subsection{Options Data}

\begin{center}
\begin{tabular}{lll}
\toprule
\textbf{Source} & \textbf{URL} & \textbf{Content} \\
\midrule
OptionMetrics & (institutional, paid) & Historical option chains \\
CBOE & \texttt{cboe.com} & Free end-of-day quotes, VIX history \\
Yahoo Finance & \texttt{finance.yahoo.com} & Basic option chains \\
\bottomrule
\end{tabular}
\end{center}

The book's option examples (Ch.~10) are self-contained: they price
from the Black-Scholes formula rather than reading market option
prices. For implied-volatility surface studies, CBOE's free VIX
history is the practical starting point.

\subsection{Microstructure and LOB Data}

\begin{center}
\begin{tabular}{lll}
\toprule
\textbf{Source} & \textbf{URL} & \textbf{Content} \\
\midrule
LOBSTER & \texttt{lobsterdata.com} & NASDAQ ITCH, 10 levels \\
Databento & \texttt{databento.com} & Paid, multi-venue L2/L3 \\
SEC EDGAR & \texttt{sec.gov/edgar} & Corporate filings (not LOB) \\
\bottomrule
\end{tabular}
\end{center}

LOBSTER is the academic standard for limit-order-book research, with
free access for academics. The book's microstructure examples
(Ch.~13) use a simulated book to keep the companion code
self-contained.

\section{Loading Data in Practice}

The \crate{quant-lib} examples all use the helpers in
\texttt{examples/common/mod.rs}, which resolve paths relative to
\texttt{CARGO\_MANIFEST\_DIR} so they work regardless of where the
binary is invoked from:

\begin{lstlisting}[language=Rust]
#[path = "common/mod.rs"]
mod common;

fn main() {
    // Synthetic single-stock CSV
    let bars = common::load_stock_csv(
        &common::stock_csv_path(env!("CARGO_MANIFEST_DIR"))
    );

    // Real B3 stock JSON
    let bars = common::load_json_ohlcv(
        &common::b3_json_path(env!("CARGO_MANIFEST_DIR"), "PETR4")
    );
}
\end{lstlisting}

To use your own data, either drop a file in
\texttt{crates/quant-lab/data/} (and it will be found by the helpers),
or write a loader that returns \texttt{Vec<OhlcvBar>} and pass it to
the example code. The \texttt{OhlcvBar} struct is:

\begin{lstlisting}[language=Rust]
pub struct OhlcvBar {
    pub date: String,
    pub open: f64,
    pub high: f64,
    pub low: f64,
    pub close: f64,
    pub volume: f64,
}
\end{lstlisting}

\section{Reproducibility and Licensing}

\begin{itemize}
    \item The B3 JSON files were downloaded from a public data provider
      and are included here for educational use; verify the provider's
      terms before redistributing.
    \item The synthetic \texttt{stock\_prices.csv} and
      \texttt{creditcard\_sample.csv} are generated by the book and are
      public domain.
    \item Public-source data (Stooq, Yahoo, Ken French, FRED) is
      generally free for research and personal use; check each
      provider's terms for commercial use.
\end{itemize}

\begin{keyinsight}
Every example in the book can be reproduced from the bundled data
alone. The public sources listed here are for readers who want to
extend the analysis to other tickers, longer windows, or alternative
asset classes. The companion code's loader abstraction means switching
data sources is a one-function change.
\end{keyinsight}